\documentclass{ctexart}
\usepackage{listings}
\usepackage{xcolor}
\usepackage{graphicx}
\usepackage{geometry}
\geometry{a4paper, margin=1in}

% 定义Verilog代码样式
\lstdefinestyle{verilog}{
    language=Verilog,
    basicstyle=\ttfamily\small,
    keywordstyle=\color{blue},
    commentstyle=\color{green!60!black},
    stringstyle=\color{red},
    numbers=left,
    numberstyle=\tiny\color{gray},
    stepnumber=1,
    numbersep=5pt,
    backgroundcolor=\color{gray!5},
    frame=single,
    rulecolor=\color{black},
    tabsize=2,
    breaklines=true,
    breakatwhitespace=false,
    captionpos=b
}

\title{8位定点CNN基本运算模块设计与功能仿真报告}
\author{}
\date{}

\begin{document}
\maketitle

\section{引言}
本报告旨在完成面向边缘端AI的8位定点CNN基本运算模块设计，包括8位无符号乘法器、8位加法器、ReLU激活模块和除以4除法器，并通过功能仿真验证各模块的正确性，为后续卷积-激活-池化全链路设计奠定基础。

\section{各模块设计与实现}

\subsection{8位无符号乘法器（卷积乘算）}
\subsubsection{功能说明}
该模块用于卷积运算中的像素值与卷积核权重相乘，输入为两个8位无符号数，输出为16位无符号乘积，以防止乘法溢出。

\subsubsection{Verilog代码实现}
\begin{lstlisting}[style=verilog, caption=mul\_8bit.v]
// mul_8bit.v
module mul_8bit (
    input  wire [7:0] a,      // 乘数A（8位无符号）
    input  wire [7:0] b,      // 乘数B（8位无符号）
    output wire [15:0] product // 乘积（16位无符号）
);
    assign product = a * b;    // 直接使用Verilog乘法运算符
endmodule
\end{lstlisting}

\subsection{8位加法器（卷积累加、池化求和）}
\subsubsection{功能说明}
该模块用于卷积结果的累加和池化运算中的求和，输入为两个8位无符号数，输出为9位无符号和，通过扩展位宽防止加法溢出。

\subsubsection{Verilog代码实现}
\begin{lstlisting}[style=verilog, caption=add\_8bit.v]
// add_8bit.v
module add_8bit (
    input  wire [7:0] a,      // 加数A（8位无符号）
    input  wire [7:0] b,      // 加数B（8位无符号）
    output wire [8:0] sum      // 和（9位无符号，防止溢出）
);
    assign sum = a + b;        // 直接使用Verilog加法运算符
endmodule
\end{lstlisting}

\subsection{ReLU激活模块（无符号数直接输出）}
\subsubsection{功能说明}
该模块用于ReLU激活函数，由于输入为8位无符号数（范围0-255），本身非负，因此无需进行负数判断，直接将输入赋值给输出即可。

\subsubsection{Verilog代码实现}
\begin{lstlisting}[style=verilog, caption=relu.v]
// relu.v
module relu (
    input  wire [7:0] in,      // 输入（8位无符号）
    output wire [7:0] out      // 输出（等于输入，因无符号数非负）
);
    assign out = in;            // 无符号数无需判断，直接输出
endmodule
\end{lstlisting}

\subsection{除以4除法器（池化平均）}
\subsubsection{功能说明}
该模块用于2×2平均池化运算，输入为10位无符号数（兼容4个8位数累加结果，最大值为4×255=1020），通过无符号数右移2位实现除以4取整，输出为8位无符号结果。

\subsubsection{Verilog代码实现}
\begin{lstlisting}[style=verilog, caption=div4.v]
// div4.v
module div4 (
    input  wire [9:0] in,      // 输入（10位，兼容4个8位数累加结果）
    output wire [7:0] out      // 输出（8位，除以4后结果）
);
    assign out = in >> 2;       // 无符号数右移2位，等价于除以4取整
endmodule
\end{lstlisting}

\section{功能仿真验证}

\subsection{仿真激励代码}
编写Testbench对四个模块进行统一测试，覆盖典型值和边界值（如0、255等）。

\begin{lstlisting}[style=verilog, caption=tb\_top.v]
// tb_top.v
`timescale 1ns/1ps  // 时间单位/精度

module tb_top;
    // 乘法器测试信号
    reg  [7:0] mul_a, mul_b;
    wire [15:0] mul_product;

    // 加法器测试信号
    reg  [7:0] add_a, add_b;
    wire [8:0] add_sum;

    // ReLU测试信号
    reg  [7:0] relu_in;
    wire [7:0] relu_out;

    // 除法器测试信号
    reg  [9:0] div4_in;
    wire [7:0] div4_out;

    // 实例化所有模块
    mul_8bit u_mul (
        .a(mul_a),
        .b(mul_b),
        .product(mul_product)
    );

    add_8bit u_add (
        .a(add_a),
        .b(add_b),
        .sum(add_sum)
    );

    relu u_relu (
        .in(relu_in),
        .out(relu_out)
    );

    div4 u_div4 (
        .in(div4_in),
        .out(div4_out)
    );

    // 测试激励
    initial begin
        // 初始化输入
        mul_a = 0; mul_b = 0;
        add_a = 0; add_b = 0;
        relu_in = 0;
        div4_in = 0;
        #10;

        // --------------------------
        // 测试1：乘法器
        // --------------------------
        mul_a = 8'd2;  mul_b = 8'd3;  #10;  // 2*3=6
        mul_a = 8'd255; mul_b = 8'd255; #10; // 255*255=65025

        // --------------------------
        // 测试2：加法器
        // --------------------------
        add_a = 8'd100; add_b = 8'd200; #10; // 100+200=300
        add_a = 8'd255; add_b = 8'd255; #10; // 255+255=510

        // --------------------------
        // 测试3：ReLU
        // --------------------------
        relu_in = 8'd123; #10; // 输入123，输出123
        relu_in = 8'd0;   #10; // 输入0，输出0

        // --------------------------
        // 测试4：除以4
        // --------------------------
        div4_in = 10'd1020; #10; // 1020/4=255
        div4_in = 10'd100;  #10; // 100/4=25

        #10;
        $finish; // 结束仿真
    end
endmodule
\end{lstlisting}
\clearpage


\begin{figure}[htbp]
    \centering
    \includegraphics[width=\textwidth]{任务1波形图.png}
    \caption{功能仿真波形图}
    \label{fig:sim_wave}
\end{figure}

\section{结论}
本报告完成了8位无符号乘法器、8位加法器、ReLU激活模块和除以4除法器的设计与实现，并通过功能仿真验证了各模块的逻辑正确性。仿真结果表明，所有模块均能正常工作，覆盖了典型值和边界值测试，可用于后续卷积-激活-池化全链路的设计与集成。

\end{document}